\section{Matching across cuffs: positions and directions}
\label{sec:matching}

The pants viewpoint suggests that orthogeodesic pieces should glue when
their boundary data match. There are two levels of matching.

First consider a cuff of length $a$, with endpoint configurations
$x_1,\ldots,x_m$ and $y_1,\ldots,y_m$ on the two sides. Choose compatible
coordinates in $\RR/a\ZZ$, reversing one induced boundary orientation
as needed. For a prescribed cyclic matching $\sigma$, a twist $\tau$
must satisfy
\begin{equation}
 y_{\sigma(j)}-x_j\equiv\tau\pmod a
       \qquad(1\le j\le m).
 \label{eq:endpointmatch}
\end{equation}
One endpoint determines $\tau$ modulo $a$; the remaining $m-1$ equations
require the cyclic spacings to agree. Integer winding data must also be
retained if the curve is marked, rather than only considered modulo a
twist.

Second, a general geodesic segment has a tangential velocity component
$p$ at an endpoint. If $S$ is the sum of segment lengths, the first
variation at a glued endpoint has the form
\begin{equation}
 \frac{\partial S}{\partial x}=p_{\rm in}-p_{\rm out}.
 \label{eq:momentummatch}
\end{equation}
Stationarity imposes tangent matching as well as position matching. For
orthogonal segments the tangential components vanish, so smooth gluing
is automatic once the endpoints match. A generic closed geodesic does
not satisfy this orthogonality condition along an arbitrarily chosen
pants decomposition.

Equivalently, boundary data live in a phase space $(s,p)$: endpoint
position and tangential component. The length of an arc is a generating
function, with
\[
 \dd L=-p_s\,\dd s+p_t\,\dd t,
 \qquad\text{and cuff gluing acts by }(s,p)\longmapsto(s+\tau,p).
\]
A closed geodesic is a periodic compatible chain of these relations.
The orthospectrum records the special $p=0$ representatives. It is
excellent input for gap identities and local geometry, but endpoint
positions and directions are needed to reconstruct arbitrary glued curves.

